\documentclass[11pt]{article}
\usepackage{fontspec}
\usepackage[french]{babel}
\usepackage{amsmath}
\usepackage{amssymb}
\usepackage{amsthm}
\usepackage{mathrsfs}
\usepackage{tikz-cd}
\usepackage{enumitem}

\newcommand{\printedpage}[1]{\marginpar{\footnotesize p.~#1}}
\newcommand{\msp}[1]{\marginpar{\footnotesize [#1]}}
\newcommand{\marginbar}{\marginpar{$\Vert$}}
\newcommand{\Hot}{\mathrm{Hot}}
\newcommand{\Hom}{\mathrm{Hom}}
\newcommand{\Cat}{\mathrm{Cat}}
\newcommand{\hot}{\mathrm{hot}}
\newcommand{\Ob}{\mathrm{Ob}}
\newcommand{\card}{\operatorname{card}}
\newcommand{\pr}{\mathrm{pr}}
\newcommand{\id}{\mathrm{id}}

\makeatletter
\renewcommand\@makefnmark{\hbox{\@textsuperscript{\normalfont(\@thefnmark)}}}
\renewcommand\@makefntext[1]{\parindent 1em\noindent\hbox{\@textsuperscript{\normalfont\@thefnmark}}\,#1}
\makeatother

\begin{document}

\printedpage{69}
% head: A. Grothendieck : Dérivateurs, Ch. XV, Édition des Universités de Montpellier II et Paris VII
% foot: Version 5 novembre 2010
où
\[
\pr_2^{I_\alpha}\ :\ I_\alpha\times X \longrightarrow X\ .
\]

\msp{page 128}Considérons maintenant le diagramme (traits pleins)
\[
\begin{tikzcd}[column sep=small]
& X\amalg X \arrow[dl, "{(\partial_0^X,\partial_1^X)}"'] \arrow[dr, "{(\partial_0^{\prime X},\partial_1^{\prime X})}"] & \\
I\times X \arrow[rr, "\varphi\times\id_X\,=\,\varphi_X"] \arrow[dr, "h"'] & & J\times X \arrow[dl, dotted, "h'"] \\
& Y\ , &
\end{tikzcd}
\]
où $\partial_0^{\prime X}$, $\partial_1^{\prime X}$ sont définis comme $\partial_0^X$, $\partial_1^X$, mais utilisant $j_0,j_1\in\Ob J$ au lieu de $i_0,i_1\in\Ob I$. S'il existait dans Cat une flèche $h'$ (pointillée) telle que $h'\circ\varphi_X=h$, il s'ensuivrait que $h'$ est une homotopie de $f_0$ et $f_1$ suivant $(J,j_0,j_1)$. Comme $J$ est 0-connexe, il s'ensuivrait aussitôt que $f_0$, $f_1$ sont homotopes au sens habituel dans Cat, donc que $f_0$ et $f_1$ ont même image dans $\Hot_W$. Mais si on ne peut trouver une flèche $h'$ dans Cat, du moins peut-on la trouver dans $\Hot_W$, puisque $\varphi_X\in W$, donc un isomorphisme dans $\Hot_W$ (et $h'$ dans $\Hom_W(J\times X,Y)$ sera même uniquement déterminé). Donc dans $\Hot_W$, on a
\[
f_0\ =\ h'\partial_0^{\prime X}\ ,\qquad\quad f_1\ =\ h'\partial_1^{\prime X}\ .
\]

\msp{page 129}Pour prouver que
\[
f_0\ =\ f_1\qquad\text{dans}\quad\Hom_W(X,Y)\ ,
\]
il suffit de voir que
\[
\partial_0^{\prime X}\ =\ \partial_1^{\prime X}\qquad\text{dans}\quad\Hom_W(X,J\times X)\ .
\]
Or en fait $\partial_0^{\prime X}$ et $\partial_1^{\prime X}$ sont homotopes, car $e\underset{j_1}{\overset{j_0}{\rightrightarrows}}J$ sont homotopes ($j_0$ et $j_1$ appartenant à la même composante connexe de $J$), donc aussi leur produit par $\id_X$.

On voit donc que pour deux objets quelconques $X$, $Y$ de Cat, toutes les flèches de $X$ dans $Y$ sont égales dans Hot. L'argument fait prouve même que toutes les flèches \textit{dans} $\Hot_W$ de $X$ dans $Y$ sont égales. Cela signifie que la catégorie $\Hot_W^*$ (où pour deux objets quelconques $X$, $Y$, $\Hom_{\Cat}(X,Y)$, et a fortiori $\Hom_W(X,Y)$ est $\neq\emptyset$) est équivalente à la catégorie ponctuelle. Quant à $\Hot_W$ elle-même, dont l'ensemble des objets est $\{\emptyset\}\amalg\Ob\Hot_W^*$, avec $\emptyset$ objet initial, elle sera donc équivalente à la catégorie ponctuelle, ou bien à $\boldsymbol{\Delta}_1$, selon que $\emptyset\longrightarrow e$ est un isomorphisme dans $\Hot_W$, ou non.

\msp{page 130}\textbf{Commentaire 9.8.} Si les axiomes Loc\,(1, 2, 3, 5)~\footnote[52]{En fait, ce commentaire vaut même en l'absence de l'hypothèse Loc\,(5), \textit{cf.} 9.9 ci-dessous.} sont satisfaits, alors l'axiome de connexité Loc\,(8) \textit{équivaut} à chacun des suivants :

\printedpage{70}
% head: A. Grothendieck : Dérivateurs, Ch. XV, Édition des Universités de Montpellier II et Paris VII
% foot: Version 5 novembre 2010
\begin{enumerate}[leftmargin=*]
\item[\textit{a})] $\Hot_W$ n'est équivalente ni à $e$, ni à $\boldsymbol{\Delta}_1$.
\item[\textit{b})] $\Hot_W$ a une infinité de classes d'isomorphie.
\item[\textit{c})] $f_0:X_0\longrightarrow Y_0$ (flèche donnée d'avance dans Cat, telle que $\pi_0(f_0)$ n'est pas un isomorphisme, et $X_0\neq\emptyset$) n'est \textit{pas} un isomorphisme de $\Hot_W$.
\item[\textit{d})] $X_0$ et $Y_0$ ne sont pas isomorphes dans $\Hot_W$ (où $X_0$, $Y_0$ sont deux objets donnés d'avance dans Cat, avec $X_0$, $Y_0$ non vides et $\card\pi_0(X_0)\neq\card\pi_0(Y_0)$).
\end{enumerate}

% marginal mark: a double vertical rule in the left margin beside the first two lines of this paragraph
\marginbar On peut donc dire que l'axiome de connexité joue le rôle d'un \textit{axiome de non trivialité} pour la catégorie $\Hot_W$ localisée de Cat. Il sera satisfait automatiquement pour tous les localiseurs fondamentaux $W$ qu'on sera amené à construire ou à introduire. Mais on n'aura pratiquement jamais à l'invoquer comme une hypothèse, en vue de certaines conclusions.

Je m'aperçois finalement que dans la proposition 9.7 l'hypothèse Loc\,(5) est inutile :

\textbf{Proposition 9.9.}~\footnote[53]{\textit{Cf.} meilleur énoncé 9.12 (p. 133 bis).} \textit{Soit $W$ un localiseur fondamental dans} Cat (i.e. \textit{on a} Loc\,(1, 2, 3) \textit{sans plus}), \textit{plus la condition que la} somme \textit{de $f_1,f_2\in W$ est dans $W$}~\footnote[54]{Cette condition Loc\,(5\,a) est vérifiée notamment si $W$ satisfait de plus l'un des deux axiomes Loc\,(4) (\textit{cf.} 9.1), ou Loc\,(6\,c$'$) (\textit{cf.} 9.3).}. \textit{Si l'axiome de connexité n'est pas satisfait, alors} $\Hot_W$ \textit{est équivalente à la catégorie ponctuelle} (\textit{si $\emptyset\longrightarrow e$ est un isomorphisme dans} $\Hot_W$~\footnote[55]{Cela signifie aussi qu'il existe $Z\neq\emptyset$ dans Cat, tel que $\emptyset\longrightarrow Z$ soit un isomorphisme dans $\Hot_W$.}) \textit{ou à $\boldsymbol{\Delta}_1$} (\textit{dans le cas contraire}), \textit{et} $\Hot_W^*$ ($\overset{\text{déf}}{=}\Hot_W\backslash\{\emptyset\}$) \textit{est équivalente à la catégorie ponctuelle.}

\msp{page 131}La condition sur les sommes assure que dans $\Hot_W$ la somme de deux objets existe, et que $\Cat\longrightarrow\Hot_W$ commute aux sommes de deux objets. Ceci dit, soit
\[
f\ :\ X\longrightarrow Y \tag*{$(*)$} % the label (*) is printed at the left margin
\]
une flèche telle que $f\in W$, mais $\pi_0(f)$ non surjectif. Supposons d'abord $X\neq\emptyset$ (cas le moins évident). Si $Y_0$ est un sommand non vide de $Y$ tel que $f$ se factorise par $Y_1=Y\backslash Y_0$, alors $f$ est somme de
\[
f_1\ :\ X\ =\ X_1\longrightarrow Y_1\qquad\quad\text{et}\qquad\quad f_0\ :\ X_0\ =\ \emptyset\longrightarrow Y_0
\]
dans Cat, donc aussi dans $\Hot_W$. Donc on a une somme $f\ =\ f_0\amalg f_1$ dans $\Hot_W$ qui est un isomorphisme, donc pour tout $Z$ dans $\Hot_W$,
\[
\Hom_W(Y_0,Z)\times\Hom_W(Y_1,Z)\ \xrightarrow{\ \sim\ }\ \underbrace{\Hom_W(X_0,Z)}_{=\,e}\times\Hom_W(X_1,Z)
\]
(\textbf{N.B.} $\emptyset$ est objet initial de $\Hot_W$, l'étant de Cat), \textit{i.e.} le composé
\[
Z(Y_0)\times Z(Y_1)\ \xrightarrow{\pr_2}\ Z(Y_1)\ \xrightarrow{Z(f_1)}\ Z(X_1)
\]
est bijectif. Comme $X\ =\ X_1\neq\emptyset$, donc aussi $Y_1\neq\emptyset$, et $Y_0\neq\emptyset$, les ensembles $Z(Y_0)$, $Z(Y_1)$, $Z(X_1)$ sont $\neq\emptyset$ si $Z\neq\emptyset_{\Cat}$. Le composé étant bijectif, la première flèche $\pr_2$ est injective, ce qui implique
\printedpage{71}
% head: A. Grothendieck : Dérivateurs, Ch. XV, Édition des Universités de Montpellier II et Paris VII
% foot: Version 5 novembre 2010
\msp{page 132}que $Z(Y_0)$ est réduit à un point (les deux facteurs $Z(Y_0)$, $Z(Y_1)$ étant $\neq\emptyset$). Cela signifie que $Y_0$ est un objet initial dans $\Hot_W^*$. Il s'ensuit notamment que $Y_0\amalg Y_0$ est aussi un objet initial, et l'application canonique $Y_0\amalg Y_0\longrightarrow Y_0$ est donc un isomorphisme dans $\Hot_W$, alors que l'application induite sur les $\pi_0$ n'est pas injective. (C'est $\pi_0(Y_0)\amalg\pi_0(Y_0)\longrightarrow\ \pi_0(Y_0)$, avec $\pi_0(Y_0)$ non vide.) Ce cas va être traité ci-dessous.

Si dans $(*)$ $X$ est vide, $Y$ non vide (puisque $\pi_0(f)$ non surjectif), alors $Y$ est objet initial de $\Hot_W$, et on voit encore que $Y\amalg Y\longrightarrow Y$ est un isomorphisme dans $\Hot_W$, bien que n'induisant pas un isomorphisme des $\pi_0$. La proposition 9.9 va donc résulter du

\textbf{Lemme 9.10.} \textit{Soit $W$ localiseur fondamental sans plus} (\textit{axiomes} Loc\,(1, 2, 3)). \textit{Supposons qu'il existe $f:I\longrightarrow J$ dans} Cat \textit{tel que} $\hot_W(f)$ \textit{soit un isomorphisme, mais que $\pi_0(f)$} ne soit pas injectif~\footnote[56]{Ou bien tel que $I$ soit 0-connexe, et $J$ non 0-connexe (\textit{i.e.} $\card\pi_0(I)\ =\ 1$, $\card\pi_0(J)\ \geq\ 2$). Ce deuxième cas se traite par la même méthode essentiellement.}. \textit{Alors} $\Hot_W^*$ \textit{est équivalente à la catégorie ponctuelle $e$, et} $\Hot_W$ \textit{est équivalente à $e$} (\textit{si $\emptyset\longrightarrow e$ est un isomorphisme dans} $\Hot_W$) \textit{ou à $\boldsymbol{\Delta}_1$} (\textit{dans le cas contraire}).

\msp{page 133}La démonstration est celle donnée pour 9.7 : on prouve que si $X$, $Y$ sont deux objets de Cat, \textit{non vides}, alors deux flèches $f_0,f_1:X\rightrightarrows Y$ sont égales dans $\Hot_W$, par un argument d'homotopie, en choisissant deux points $i_0$, $i_1$ de $I$ qui sont dans deux composantes connexes distinctes, et tels que leurs images $j_0$, $j_1$ dans $J$ soient dans la même composante.

\textbf{Proposition 9.11.} \textit{Soit $W$ un localiseur fondamental} (i.e. \textit{on a} Loc\,(1, 2, 3)). \textit{Alors les conditions suivantes sont équivalentes.}
\begin{enumerate}[leftmargin=*]
\item[a)] \textit{On a} Loc\,(5\,bis) (\textit{axiome de décomposition, forme forte}) \textit{et} Loc\,(8) (\textit{axiome de connexité}).
\item[b)] \textit{On a} Loc\,(5\,bis), \textit{et} $\Hot_W$ \textit{n'est pas équivalente à $e$ ou à $\boldsymbol{\Delta}_1$} (cf. \textit{commentaire} 9.8).
\item[c)] \textit{Pour toute $f:X\longrightarrow Y$, on a $f\in W$ si et seulement si $\pi_0(f)$ est un isomorphisme, et pour toute composante connexe $X_i$ de $X$, si $Y_i$ désigne la composante connexe correspondante de $Y$, la flèche induite $f_i:X_i\longrightarrow Y_i$ est dans $W$.}
\end{enumerate}

La raison pour laquelle je n'isole pas cette propriété en un des axiomes Loc\,(\ldots), c'est que ce n'est \textit{pas} un axiome de stabilité, mais un axiome de nature mixte.

\msp{page 133 bis}\textbf{Proposition 9.12.} \textit{Soit $W$ un localiseur fondamental} (\textit{sans plus}), i.e. \textit{satisfaisant} Loc\,(1, 2, 3). \textit{Les conditions suivantes sont équivalentes.}
\begin{enumerate}[leftmargin=*]
\item[a)] \textit{$W$ satisfait l'axiome de connexité} Loc\,(8), i.e. $f\in W\Longrightarrow\pi_0(f)$ \textit{bijectif.}
\item[b)] \textit{Si $f,g:X\rightrightarrows Y$ dans} Cat \textit{sont telles que} $\hot_W(f)=\hot_W(g)$, \textit{alors} $\pi_0(f)=\pi_0(g)$.
\item[c)] \textit{Il existe $X$, $Y$ dans} Cat \textit{tel que l'application} $\hot_W^{X,Y}:\Hom_{\Cat}(X,Y)\longrightarrow\Hom_W(X,Y)$ \textit{ne soit pas constante.}
\item[c$'$)] \textit{La sous-catégorie} $\Hot_W^*$ \textit{de} $\Hot_W$, \textit{formée des objets non vides, n'est pas équivalente à la catégorie ponctuelle.}
\item[c$''$)] \textit{La catégorie} $\Hot_W$ \textit{n'est pas équivalente à la catégorie $e$ ni à la catégorie $\boldsymbol{\Delta}_1$.}
\item[d)] \textit{L'ensemble des classes d'isomorphie d'objets de} $\Hot_W$ \textit{est infini.}
\end{enumerate}

\printedpage{72}
% head: A. Grothendieck : Dérivateurs, Ch. XV, Édition des Universités de Montpellier II et Paris VII
% foot: Version 5 novembre 2010
\textsc{Démonstration}. On a $(a)\Longrightarrow(b)$, puisque $(a)$ équivaut à dire que le foncteur $\pi_0:\Cat\longrightarrow\mathrm{Ens}$ se factorise par $\Hot_W$. D'ailleurs le foncteur correspondant
\[
\Hot_W\ \longrightarrow\ (\mathrm{Ens})
\]
est essentiellement surjectif -- il y a un inverse à isomorphisme près, $I\longmapsto\hot_W(I)$. Cela prouve que $(a)$ implique $(d)$, et on a $(d)\Longrightarrow(c')$, $(d)\Longrightarrow(c'')$ tautologiquement. Un argument plus ou moins tautologique montre d'ailleurs que si $\Hot_W^*$ est équivalente à $e$, alors $\Hot_W$ est équivalente à $e$ ou à $\boldsymbol{\Delta}_1$ (car $\emptyset$ est objet initial de $\Hot_W$) -- donc $(c'')\Longrightarrow(c')$. Donc on a plus ou moins tautologiquement
\[
\begin{tikzcd}
(\mathrm{a}) \arrow[r, Rightarrow] \arrow[d, Rightarrow] & (\mathrm{d}) \arrow[r, Rightarrow] & (\mathrm{c}'') \arrow[r, Rightarrow] & (\mathrm{c}') \\
(\mathrm{b}) \arrow[rr, Rightarrow] & & (\mathrm{c}) \arrow[ru, Rightarrow] &
\end{tikzcd}
\]
où les implications $(b)\Longrightarrow(c)\Longrightarrow(c')$ sont elles aussi triviales \texttt{[c'est plutôt l'implication (c$'$) $\Longrightarrow$ (c) qui est triviale~; l'implication (c) $\Longrightarrow$ (c$'$) est démontrée après l'énoncé du corollaire 9.13, avant la preuve de ce dernier]}. Donc
\msp{page 133 ter}tout revient à prouver que $(c)\Longrightarrow(a)$, ou encore :

\textbf{Corollaire 9.13.} \textit{Si l'axiome de connexité n'est pas satisfait, alors si $X$, $Y$ sont deux objets de} Cat, \textit{toutes les flèches $X\longrightarrow Y$ dans} Cat \textit{sont égales dans} $\Hot_W$, i.e. \textit{l'application} $\Hom_{\Cat}(X,Y)\longrightarrow\Hom_W(X,Y)$ \textit{est constante.}

Cela implique donc que tous les objets de Cat$^*$, \textit{i.e.} tous les objets non vides de Cat, sont isomorphes dans $\Hot_W$, et on trouve un isomorphisme canonique entre elles, avec transitivité. Soit alors $\mathcal{H}$ la catégorie quotient de Cat, ayant mêmes objets que Cat, et ayant exactement une flèche entre deux objets non vides de Cat et exactement une flèche de $\emptyset$ dans tout autre objet de Cat (et pas de flèche d'un objet non vide de Cat dans $\emptyset$). C'est donc une catégorie équivalente à $\boldsymbol{\Delta}_1$. On a vu que $\Cat\longrightarrow\Hot_W$ se factorise en
\[
\Cat\ \xrightarrow{\text{can.}}\ \mathcal{H}\ \longrightarrow\ \Hot_W\ .
\]
Si $\Cat\longrightarrow\mathcal{H}$ transforme toute flèche appartenant à $W$ en isomorphisme, alors il est immédiat que $\mathcal{H}$ est solution au problème universel définissant $\Hot_W$, donc que $\mathcal{H}\longrightarrow\Hot_W$ est un isomorphisme. Cela signifie aussi que $W$ ne contient pas de flèche $\emptyset\longrightarrow X$ avec $X$ non vide (auquel cas $W$ contient toutes les flèches). Donc dans ce cas $\Hot_W\simeq\boldsymbol{\Delta}_1$. Dans le cas contraire, soit $\mathcal{H}'$ la catégorie déduite de $\mathcal{H}$ en inversant les flèches $\emptyset\longrightarrow X$, \textit{i.e.} le groupoïde 1-connexe sur l'ensemble d'objets $\Ob\Cat$. Ce n'est plus une catégorie quotient de Cat, mais on a encore une factorisation
\[
\Cat\ \xrightarrow{\text{épi}}\ \mathcal{H}\ \ \xrightarrow{\text{loc.}}\ \mathcal{H}'\ \longrightarrow\ \Hot_W\ ,
\]
\printedpage{73}
% head: A. Grothendieck : Dérivateurs, Ch. XV, Édition des Universités de Montpellier II et Paris VII
% foot: Version 5 novembre 2010
\msp{page 133 quater}et cette fois c'est $\mathcal{H}'$ qui est solution du problème universel ($\Cat\longrightarrow\mathcal{H}'$ transformant flèches de $W$ en isomorphismes). Donc dans ce cas, $\mathcal{H}'\longrightarrow\Hot_W$ est un isomorphisme, et $\Hot_W$ est équivalente à $e$. Dans les deux cas, $\Hot_W^*$ est équivalente à $e$.

Donc, la proposition 9.12 résulte du corollaire 9.13, qu'il s'agit à présent de démontrer.

Soit $f:I\longrightarrow J$ dans Cat telle que $\pi_0(f)$ ne soit pas injective. On a vu dans la démonstration de 9.7 comment on en déduit (en choisissant deux éléments $i_0$, $i_1$ dans deux composantes connexes distinctes de $I$, telles que $j_0=f(i_0)$ et $j_1=f(i_1)$ soient dans une même composante connexe de $J$) que pour $X$, $Y$ dans Cat, deux flèches $f_0,f_1:X\rightrightarrows Y$ sont toujours égales dans $\Hot_W$ (\textit{cf.} lemme 9.10).

Soit maintenant
\[
f\ :\ I\longrightarrow J
\]
dans Cat telle que $\pi_0(f)$ ne soit pas surjective, donc $J=J_0\amalg J_1$, avec $f$ se factorisant par $J_0$, et $J_1\neq\emptyset$. Considérons alors la catégorie discrète à deux objets,
\[
\varepsilon\ =\ \{0,1\}\ ,
\]
et deux flèches
\[
g_0,\ g_1\ :\ J\ \rightrightarrows\ \varepsilon\ ,
\]
$g_0$ étant constante de valeur 0, et $g_1$ étant égale à 0 sur $J_0$, à 1 sur $J_1$. On a
\[
g_0f\ =\ g_1f\ ,
\]
donc, comme $f\in W$,
\[
g_0\ =\ g_1\qquad\quad\text{dans}\quad\Hot_W\ .
\]
On en conclut, en prenant un point $j_1$ dans $J_1$, et l'inclusion $\alpha_{j_1}:e\longrightarrow J$ correspondante, que
\msp{page 133 quinquies}
\[
g_0\alpha_{j_1}\ =\ g_1\alpha_{j_1}\ ,
\]
\textit{i.e.} les deux sections de $\varepsilon$ sur $e$ sont égales dans $\Hot_W$. Comme le foncteur $\Cat\longrightarrow\Hot_W$ commute aux produits finis \texttt{[cette propriété n'est pas établie sous les seules conditions Loc\,(1,2,3)~; elle est démontrée dans Pursuing Stacks, section 64, en remplaçant Loc\,(3) par l'axiome plus fort Loc\,(3\,bis)]}, il s'ensuit que pour $X$ dans Cat, les deux flèches
\[
X\ \underset{\partial_1^X}{\overset{\partial_0^X}{\rightrightarrows}}\ X\times\varepsilon\ =\ X\amalg X
\]
sont égales dans $\Hot_W$. Si alors
\[
f_0,\ f_1\ :\ X\ \rightrightarrows\ Y
\]
sont deux flèches dans Cat, elles sont de la forme
\[
f_0\ =\ h\partial_0^X\ ,\qquad\quad f_1\ =\ h\partial_1^X\ ,
\]
avec
\[
h\ =\ (f_0,f_1)\ :\ X\amalg X\ \longrightarrow\ X\ , % [sic] X\amalg X \longrightarrow X -> X\amalg X \longrightarrow Y, h\partial_i^X = f_i : X \to Y requires target Y
\]
et comme $\partial_0^X=\partial_1^X$ dans $\Hot_W$, on en conclut $f_0=f_1$ dans $\Hot_W$, \quad \textit{q.e.d.}

\end{document}
